\documentclass[11pt,letterpaper]{article}
\usepackage[margin=0.85in]{geometry}
\usepackage{amsmath,amssymb,graphicx,booktabs,float}
\usepackage[font=small,labelfont=bf]{caption}
\usepackage[hidelinks]{hyperref}
\setlength{\parindent}{0pt}
\setlength{\parskip}{4pt}
\setlength{\textfloatsep}{10pt}
\setlength{\intextsep}{8pt}
\raggedbottom
\input{report_experiments/report_numbers.tex}

\begin{document}
\begin{center}
    {\Large\bfseries CSE 256 PA1: Sentiment Classification and Word Representations}\\[5pt]
    Jiaqi Wang \quad $\vert$ \quad October 23, 2026
\end{center}

\section{Task and model}
I study binary sentiment classification on the supplied movie-review snippets from
the Stanford Sentiment Treebank~\cite{socher}. The provided split contains 6,920
training examples and 872 development examples. The supplied reader lowercases
the tokenized sentences. Development labels are used for model assessment, not
gradient updates; no held-out test results are reported.

My deep averaging network (DAN) follows the averaging-and-classification idea
of Iyyer et al.~\cite{iyyer}. For token IDs $t_1,\ldots,t_L$, let $m_i$ be zero
for padding and one otherwise. The sentence representation and class scores are
\[
 h=\frac{\sum_{i=1}^{L}m_i e(t_i)}{\max(1,\sum_{i=1}^{L}m_i)},
 \qquad
 z=W_2\,\operatorname{Dropout}\!\left(\operatorname{ReLU}(W_1h+b_1)\right)+b_2.
\]
All DANs use 300-dimensional embeddings, one 100-unit hidden layer, dropout 0.2,
and two output scores. Adam uses learning rate $10^{-3}$ with PyTorch's other
default settings. Cross-entropy consumes the raw scores; prediction is
$\arg\max_k z_k$. Every five epochs, loss and accuracy are measured on both
splits with dropout disabled and without gradients. ``Best dev'' means the
maximum over these recorded evaluations, not an independently estimated test
accuracy or a saved model checkpoint.
 
\section{Part 1: GloVe versus learned word embeddings}
\textbf{Setup.} The first model uses the supplied 300-dimensional GloVe
vectors~\cite{glove}. Embeddings are frozen, so sentence means are precomputed
and only the classifier is trained. Missing words use the supplied zero UNK
vector. The second model uses the same 14,923-entry word index, including PAD
and UNK, but initializes its embedding matrix randomly and trains it jointly
with the classifier. Only the index is reused; no GloVe vector values initialize
this model. PAD is excluded from averaging. Both original experiments run for
200 epochs with one full-batch parameter update per epoch.

\begin{table}[H]
\centering\small
\begin{tabular}{lrrrr}
\toprule
Model & Final train (\%) & Final dev (\%) & Best dev (\%) & Best epoch\\
\midrule
\PartOneRows
\bottomrule
\end{tabular}
\caption{Original full-batch runs. All final values are from epoch 200.}
\label{tab:part1}
\end{table}

\textbf{Results.} Frozen GloVe reaches \GFullBest\% best development accuracy,
exceeding the assignment's 77\% target. Random embeddings reach \RFullBest\%.
The learning curves in Figure~\ref{fig:words} distinguish fitting from
generalization: the random model's final training accuracy is \RFullTrain\%,
yet its final development accuracy is only \RFullFinal\%. Its development loss
is lowest at epoch 55 (0.5370), then rises to 1.1575 while training loss falls
to 0.0031. Development accuracy can remain similar while cross-entropy worsens,
because cross-entropy also penalizes the confidence of incorrect predictions.

\begin{figure}[H]
\centering
\includegraphics[width=\linewidth]{report_experiments/figures/word_curves.png}
\caption{Original word-level runs: training and development cross-entropy
(left) and accuracy (right). Solid lines denote train; dashed lines denote dev.}
\label{fig:words}
\end{figure}

\textbf{Interpretation.} Pretrained representations provide a useful starting
point for this small labeled dataset. The frozen model trains only 30,302
classifier parameters, whereas the random model trains 4,507,202 parameters,
including embeddings. Its faster reduction in training loss therefore does not
show better generalization. This comparison changes both initialization and
embedding trainability; it does not isolate initialization alone. The original
runs did not record a random seed or training wall time. They are reported as
single runs; controlled, timed experiments follow below.

\section{Part 2: BPE and vocabulary size}
\textbf{Tokenizer.} I implement character-based byte pair encoding (BPE),
following the subword approach of Sennrich et al.~\cite{bpe}. Each distinct
training word begins as characters plus an end-of-word symbol \texttt{</w>}.
At each iteration I count adjacent symbol pairs, weighted by word frequency,
select the most frequent pair (lexicographic tie-breaking), and merge all its
non-overlapping occurrences within words. The ordered rules and vocabulary are
fitted on training text only. Both splits are encoded by applying those rules
in order, without crossing word boundaries. Unknown characters map to UNK.
Earlier symbols are retained, so reported vocabulary sizes count all allocated
symbols, including intermediate merges, PAD, UNK, and the word-end marker.

\textbf{Controlled protocol.} I compare target vocabularies 1,000, 3,000, and
5,000 with word-level DANs. Every run uses seed 42, batch size 64, 30 epochs,
the shared classifier above, and two PyTorch CPU threads. Thus each epoch has
109 parameter updates, unlike the original full-batch runs. All subword
embeddings are random and trainable. The runs use Windows 11, Python
\PythonVersion, and PyTorch \TorchVersion. These are single-seed comparisons,
not estimates of statistical significance.

\begin{table}[H]
\centering\small
\begin{tabular}{lrrrrr}
\toprule
Representation & Vocabulary & Best dev (\%) & Best epoch & Final dev (\%) & DAN run (s)\\
\midrule
\ControlledRows
\bottomrule
\end{tabular}
\caption{Matched 30-epoch, batch-64 experiments. Word models use the supplied
word index; BPE learns its vocabulary from train. DAN time includes periodic
evaluation and excludes data preparation.}
\label{tab:controlled}
\end{table}

\begin{table}[H]
\centering\small
\begin{tabular}{rrrrrr}
\toprule
Vocabulary & Merges & Mean train length & BPE fit (s) & Total prep (s) & Trainable params\\
\midrule
\EfficiencyRows
\bottomrule
\end{tabular}
\caption{BPE efficiency on the same machine. Length counts non-PAD subwords.
Total preparation includes BPE fit, encoding, and model setup.
Parameters include the embedding matrix and classifier.}
\label{tab:efficiency}
\end{table}

\textbf{Results and tradeoffs.} The highest observed BPE development accuracy is
\BestBPEAccuracy\%, attained at vocabulary sizes of \BestBPEVocab\ symbols. The matched
random word model reaches \ControlledRandomBest\%, and frozen GloVe reaches
\ControlledGloveBest\%. Moving from 1,000 to 3,000 symbols improves best accuracy
by 1.95 percentage points; 5,000 gives no further peak gain despite higher
preparation cost. BPE does not outperform the word-level models in these runs.
Increasing BPE vocabulary from 1,000 to 5,000 reduces
mean training sequence length from \SmallBPELength\ to \LargeBPELength\,
but increases the embedding matrix and the number of tokenizer merges.
Smaller units share parameters across related words, whereas more complete-word
units can preserve word-specific sentiment. These are plausible mechanisms;
the experiments measure the resulting accuracy, not the mechanisms separately.
Subword averaging also weights words with more pieces more heavily, so it is
not equivalent to averaging one vector per original word.

\begin{figure}[H]
\centering
\includegraphics[width=\linewidth]{report_experiments/figures/bpe_curves.png}
\caption{Development curves for the three BPE vocabulary sizes under the
matched protocol. Best accuracy and smallest loss need not occur at the same epoch.}
\label{fig:bpe}
\end{figure}

The tokenizer implementation recounts pairs and scans word representations
after every merge. Its fit cost consequently grows with the merge budget.
Encoding repeated words is cached. Shorter encoded sequences do not guarantee
lower end-to-end cost, because vocabulary growth increases both tokenizer work
and embedding parameters. Table~\ref{tab:efficiency} reports those costs
separately. Development monitoring would support early stopping; the reported
runs themselves always complete their specified epoch budgets.

\section{Part 3: Understanding Skip-Gram}
Let $v_x$ be a center-word vector and $c_y$ a context vector. The model and
training objective are
\[
 P(y\mid x)=\frac{\exp(v_x^\top c_y)}{\sum_{y'\in V}\exp(v_x^\top c_{y'})},
 \qquad \mathcal{L}=\sum_{(x,y)\in D}\log P(y\mid x).
\]
The vocabulary in both questions is $V=\{\text{a},\text{the},\text{dog},\text{cat}\}$,
and a window of one produces contexts on both sides within each sentence.

\subsection*{3a. Maximum-likelihood probabilities for \texttt{the}}
For the three sentences ``the dog'', ``the cat'', and ``a dog'', the examples
whose center is \texttt{the} are $(\text{the},\text{dog})$ and
$(\text{the},\text{cat})$. Its objective is
$\log p_{\mathrm{dog}}+\log p_{\mathrm{cat}}$.
Since $p_{\mathrm{dog}}+p_{\mathrm{cat}}\leq 1$, their product is at most
$((p_{\mathrm{dog}}+p_{\mathrm{cat}})/2)^2\leq 1/4$.
Equality gives the empirical conditional distribution:
\[
 P^*(\text{dog}\mid\text{the})=P^*(\text{cat}\mid\text{the})=\tfrac12,
 \qquad P^*(\text{a}\mid\text{the})=P^*(\text{the}\mid\text{the})=0.
\]
These probabilities maximize likelihood without imposing a particular
embedding parameterization.

\subsection*{3b. A nearly optimal center vector with fixed contexts}
The specified context vectors are
$c_{\mathrm{dog}}=c_{\mathrm{cat}}=(0,1)$ and
$c_{\mathrm{a}}=c_{\mathrm{the}}=(1,0)$.
Choose $v_{\mathrm{the}}=(0,5)$. Dog and cat each receive score 5; a and the
each receive score 0. Consequently,
\[
 q=\frac{e^5}{2e^5+2}=0.4966536,\qquad
 r=\frac{1}{2e^5+2}=0.0033464.
\]
The two desired contexts receive $q$ each and the others receive $r$ each.
Every probability differs from its optimum by only $0.0033464<0.01$.
More generally, $v_{\mathrm{the}}=(0,M)$ approaches the optimum as
$M\rightarrow\infty$. Finite softmax scores cannot assign exact zero
probability to the unobserved contexts.

\subsection*{3c. Training examples for the four-sentence corpus}
Adding ``a cat'' yields eight directed training examples, each once:
\[
\begin{array}{ll}
(\text{the},\text{dog}), & (\text{dog},\text{the}),\\
(\text{the},\text{cat}), & (\text{cat},\text{the}),\\
(\text{a},\text{dog}), & (\text{dog},\text{a}),\\
(\text{a},\text{cat}), & (\text{cat},\text{a}).
\end{array}
\]
Each center has two observations, one for each word in the opposite group:
$A=\{\text{a},\text{the}\}$ and $B=\{\text{dog},\text{cat}\}$.
Thus its optimal conditional probabilities are $1/2$ for each opposite-group
word and zero for each same-group word.

\subsection*{3d. Word and context vectors approaching the global optimum}
Use separate center and context vectors as follows:
\begin{center}
\begin{tabular}{lcc}
\toprule
Words & Center vector $v$ & Context vector $c$\\
\midrule
\texttt{a}, \texttt{the} & $(0,5)$ & $(1,0)$\\
\texttt{dog}, \texttt{cat} & $(5,0)$ & $(0,1)$\\
\bottomrule
\end{tabular}
\end{center}
For every center, opposite-group contexts have score 5 and same-group contexts
have score 0. Hence the probabilities are again $q$ and $r$ from 3b, respectively,
and all are within 0.01 of the empirical optima. The joint construction works
for all four centers simultaneously. Its summed log-likelihood is
$8\log q\approx-5.5989002$, compared with the supremum
$8\log(1/2)\approx-5.5451774$. Replacing 5 by $M$ and taking
$M\rightarrow\infty$ approaches that supremum.

\begin{thebibliography}{4}
\small
\bibitem{iyyer} M. Iyyer, V. Manjunatha, J. Boyd-Graber, and H. Daum\'e III.
Deep Unordered Composition Rivals Syntactic Methods for Text Classification.
ACL, 2015. \url{https://aclanthology.org/P15-1162/}
\bibitem{glove} J. Pennington, R. Socher, and C. D. Manning.
GloVe: Global Vectors for Word Representation. EMNLP, 2014.
\url{https://aclanthology.org/D14-1162/}
\bibitem{bpe} R. Sennrich, B. Haddow, and A. Birch.
Neural Machine Translation of Rare Words with Subword Units. ACL, 2016.
\url{https://aclanthology.org/P16-1162/}
\bibitem{socher} R. Socher et al.
Recursive Deep Models for Semantic Compositionality Over a Sentiment Treebank.
EMNLP, 2013. \url{https://aclanthology.org/D13-1170/}
\end{thebibliography}
\end{document}
